% ECONOMICS_PROBLEM: {"id": "D02-416", "title": "Predicting institutional turning points", "jel_code": "D02", "source_id": "OP-0380"}
% !TeX program = xelatex
\documentclass[11pt,a4paper]{article}
\usepackage[margin=23mm]{geometry}
\usepackage{fontspec}
\setmainfont{Latin Modern Roman}
\usepackage{amsmath,amssymb,mathtools}
\usepackage{xcolor,enumitem,booktabs,longtable,array,xurl,hyperref}
\definecolor{muted}{HTML}{53636C}
\hypersetup{colorlinks=true,urlcolor=blue,linkcolor=blue}
\setlength{\parindent}{0pt}
\setlength{\parskip}{5pt}
\newcommand{\R}{\mathbb R}
\newcommand{\E}{\mathbb E}
\newcommand{\N}{\mathbb N}
\newcommand{\ind}{\mathbf 1}
\newcommand{\Var}{\operatorname{Var}}
\newcommand{\Cov}{\operatorname{Cov}}
\newcommand{\supp}{\operatorname{supp}}
\DeclareMathOperator*{\argmin}{arg\,min}
\DeclareMathOperator*{\argmax}{arg\,max}
\newcommand{\entrymeta}[3]{{\sffamily\small\color{muted}\textbf{\texttt{#1}}\quad|\quad #2\par #3\par}}
\newcommand{\Problemref}[1]{\texttt{#1}}
\newcommand{\JELref}[2]{\texttt{#2} (JEL \texttt{#1})}
\begin{document}
\section*{Predicting institutional turning points}
\textbf{Problem ID:} \texttt{D02-416}\quad \textbf{JEL:} \texttt{D02}\par
% BEGIN ECONOMICS STATEMENT
\label{problem:OP-0380}
{\sffamily\small\color{muted}Original subject group: Political economy and institutions\par}
\label{definitions:problem:30007662}
\textbf{Audit classification:} Published research agenda. Working paper and published metadata checked. Real-time belief measurement is an explicit next step; entropy, calibration and causal intervention definitions are editorial.
\subsection*{Definitions and precise research target}
\textbf{Complete editorial problem.} Fix a panel of polities with institutional states \(R_{ct}\) measured by a declared rule at date \(t\). Citizen \(i\) in polity \(c\) reports a probability distribution \(b_{ict}(r)\) over possible institutions \(r\) at horizon \(T\). Define a prespecified uncertainty index \(H_{ct}\) using within-person entropy and between-person disagreement, and let \(X_{ct}\) contain information observed before outcomes. The prediction task is to estimate \(\Pr(R_{cT}\ne R_{ct}\mid H_{ct},X_{ct})\) and test calibration in episodes selected before institutional changes are known. A critical juncture may exist even when the status quo survives, so samples cannot contain only successful transitions. For a feasible binary information or coordination intervention \(z\in\{0,1\}\), define the causal estimand \(\tau_R=E[1\{R_{cT}(1)\in\mathcal I\}-1\{R_{cT}(0)\in\mathcal I\}]\), where \(\mathcal I\) is a declared set of inclusive institutional states. Distinguish predictive uncertainty from uncertainty causing collective action. Measure higher-order beliefs where relevant, address survey selection and repression, and predefine durability. The cited authors propose real-time belief measurement; this quantified forecasting and intervention problem adapts their agenda and does not claim a universal sufficient statistic for institutional change.

\textbf{Definitions and admissible scope.} Let the institutional state space \(\mathcal R\) be finite and every elicited \(b_i\) sum to one. Set \(\bar b=\sum_i\mu_i b_i\), \(H(b)=-\sum_r b(r)\log b(r)\) with \(0\log0=0\), within-person uncertainty \(H_W=\sum_i\mu_iH(b_i)\), and disagreement \(H_B=H(\bar b)-H_W\ge0\). These are distinct predictors; neither is defined as actual future regime change. A forecast \(\widehat p\) is calibrated if \(P(R_T\ne R_t\mid\widehat p)=\widehat p\) in a declared prospective evaluation population. Institutional inclusion is a prespecified subset \(\mathcal I\subseteq\mathcal R\), measured independently of the forecast. A causal information or coordination intervention is a feasible assigned policy with its own potential-outcome law, not an intervention setting entropy directly. For each polity and date use normalized citizen weights \(\mu_{ict}\), define the preceding entropy quantities within that \((c,t)\), and set the predictor \(H_{ct}=(H_{W,ct},H_{B,ct})\) rather than an unspecified scalar combination. Use a fixed finite \(\mathcal R\), a prespecified inclusive subset \(\mathcal I\), and a fixed future date \(T>t\). The episode-sampling, missing-response, forecast-training, and held-out evaluation laws are inputs. Calibration is an out-of-sample conditional-frequency property under that population. A durable inclusive outcome may be \(1\{R_{cs}\in\mathcal I\text{ for every }s=T,\ldots,T+h\}\), with integer \(h\ge0\) fixed in advance, as a separate extension of \(\tau_R\).
\subsection*{Variants and assumption changes}
\begin{enumerate}
\item Real-time measurement (source-motivated): repeatedly elicit distributions before outcomes, sample episodes without conditioning on successful regime change, and test predictive value of \((H_W,H_B)\) relative to other known predictors.
\item Coordination and durability (editorial): randomize a feasible credible-information intervention and measure participation, higher-order beliefs, regime state, and survival over a fixed horizon. Separate calibration gains from evidence that beliefs cause institutional change.
\end{enumerate}
\subsection*{References and source audit}
\textbf{Support and access.} Working paper and published metadata checked. Real-time belief measurement is an explicit next step; entropy, calibration and causal intervention definitions are editorial. \textbf{Locator:} CESifo DP 10833 (2023), Section 3.1 p. 7; published 2024 DOI separately checked. \textbf{Versions:} Published 2024 in Annual Review of Economics 16, 105--131, DOI 10.1146/annurev-economics-091823-031317; this DOI is not a 2023 publication. 
\begin{enumerate}\raggedright
\item\hypertarget{definitions:source:30007662:1}{} Michael Callen; Jonathan L. Weigel; Noam Yuchtman. 2023. \emph{Experiments about Institutions}. CESifo Working Paper 10833, December 2023, Section 3.1, p. 7.
\url{https://www.ifo.de/DocDL/cesifo1_wp10833.pdf}.
\item\hypertarget{definitions:source:30007662:2}{} Michael Callen; Jonathan L. Weigel; Noam Yuchtman. 2024. \emph{Experiments About Institutions}. Published citation and abstract, Annual Review of Economics 16, 105--131.
\url{https://www.annualreviews.org/content/journals/10.1146/annurev-economics-091823-031317}.
DOI: \url{https://doi.org/10.1146/annurev-economics-091823-031317}.
\end{enumerate}

\subsection*{Common conventions: Source mathematical conventions}

These conventions apply to each entry unless explicitly replaced.

\textbf{Models and identification.} A primitive model \(m\) specifies all probability laws, feasible choices, information, equilibrium and measurement rules used by its target. The admissible class \(\mathcal M\) is fixed independently of outcomes. If \(P_O(m)\) is its observable probability law and \(\tau(m)\) the target, its identified set at observable law \(P\) is
\[
 \mathcal I_\tau(P)=\{\tau(m):m\in\mathcal M,\ P_O(m)=P\}.
\]
Point identification means that this set is a singleton. An empty set signals incompatibility of data and assumptions. Coordinate bounds need not describe the joint set. Bounds are sharp only if every retained target is attainable or arbitrarily approachable by compatible models. Finite bounds require explicit restrictions on unobserved quantities; unrestricted confounding, hidden wealth, extrapolation or arbitrary equilibrium selection can yield uninformative sets.

\textbf{Causal interventions.} A potential outcome \(Y_i(p)\) is the outcome under a complete policy regime \(p\), including timing, financing, information and market spillovers. The target population and units are fixed before treatment. With interference, use the complete assignment vector or a justified exposure mapping; individual no-interference notation alone cannot identify an economy-wide intervention. A difference of mean potential outcomes does not require the cross-world joint distribution. Conditioning on post-treatment migration, survival, enrollment or employment changes the estimand and needs separate justification.

\textbf{Statistical statements.} An estimator, confidence set, forecast or test specifies a sampling law, model class and loss/coverage criterion. Uniform \((1-\alpha)\)-coverage means \(\inf_{m\in\mathcal M}\Pr_m\{\tau(m)\in C_n\}\geq1-\alpha\), or an explicitly asymptotic version. The whole parameter space trivially has coverage, so informativeness requires a stated width, risk or power objective. A forecasting improvement is relative to a named benchmark, information set and evaluation population.

\textbf{Equilibrium and optimization.} An equilibrium comprises feasible plans, optimality, beliefs where needed, budget constraints and all market/resource clearing. The primitives or a stated benchmark define each correspondence \(\mathcal E(p;m)\). Nonemptiness is assumed or proved, never inferred just from compact policy sets. A selected-equilibrium target uses a declared selection rule; a robust target quantifies over all permitted equilibria. Use suprema or infima unless attainment is established. An \(\varepsilon\)-optimal policy is feasible and within \(\varepsilon\) of the finite optimum value. Report an empty feasible set separately.

\textbf{Welfare and accounting.} Normative utility, interpersonal normalization, social weights, numeraire, discounting and the use of public receipts are primitives. Behavioral utility need not coincide with the welfare criterion. Household welfare includes ownership income and rents once; adding profits again to compensating variation generally double-counts them. A monetary equivalent is defined by an explicit transfer and price convention, with monotonicity and range conditions. Average effects, mechanism contributions, transfers and welfare losses are different targets. Infinite sums require convergence and dynamic feasible plans require terminal or no-Ponzi conditions.

\textbf{Source versions.} A working-paper date, revision date and journal publication year are distinguished. A DOI for a journal version is not automatically the DOI of an earlier manuscript. Locators refer to the stated version and printed pages where specified. A metadata-only check does not verify a claim about a full-text conclusion. Every entry's source subsection narrows the provenance claim accordingly.
% END ECONOMICS STATEMENT
\end{document}
